\documentclass{article}
\counterwithin{equation}{section}
\newcounter{example}
\counterwithin*{example}{section}
\newenvironment{example}[1][]{%
	\stepcounter{example}%
	\par\vspace{5pt}\noindent
	\fbox{\textbf{Example~\thesection.\theexample}}%
	\hrulefill\par\vspace{10pt}\noindent\rmfamily}%
{\par\noindent\hrulefill\vrule width10pt height2pt depth2pt\par}
\begin{document}
	\section{First equation}
	Some introductory text...
	\begin{example}
		\begin{equation}
			f(x)=\frac{x}{1+x^2}
		\end{equation}
	\end{example}
	
	\subsection{More detail}
	\begin{example}
		Here we discuss
		\begin{equation}
			f(x)=\frac{x+1}{x-1}
		\end{equation}
		... but don't say anything
	\end{example}
	
	\subsubsection{Even more detail}
	\begin{example}
		\begin{equation}
			f(x)=\frac{x^2}{1+x^3}
		\end{equation}
		\begin{equation}
			f(x+\delta x)=\frac{(x+\delta x)^2}{1+(x+\delta x)^3}
		\end{equation}
		\begin{equation}
			f(x+\delta x)=\frac{(x+\delta x)^0}{1+(x+\delta x)^2}
		\end{equation}
		.. . Say something Pet
	\end{example}
	
	\section{Third equation}
	\begin{example}
		The following function...
		\begin{equation}
			f_1(x)=\frac{x+1}{x-1}
		\end{equation}
		..is a function
		\begin{equation}
			f_2(x)=\frac{x^2}{1+x^3}
		\end{equation}
		\begin{equation}
			f_3(x)=\frac{3+ x^2}{1-x^3}
		\end{equation}
	\end{example}
\end{document}